\begin{figure}[t]\centering\includegraphics[width=\columnwidth]{sweeps.pdf}
\caption{Left: FNO test error versus retained Fourier modes. Right: test error versus training-set size for FNO and DeepONet (3 seeds, mean $\pm$ std, log scales).}\label{fig:sweeps}\end{figure}

\begin{table*}[t]\centering\caption{Sensitivity and ablation studies, seeds 0--2, mean $\pm$ std of relative $L_2$ at 128 points. Default rows are the main-group runs for those seeds. ``train'' is the training error of the restored best-validation checkpoint. ``change'' is the relative change of the mean test error against the default row of the same model, and Welch $p$ is the unpaired test against that row (three values each).}\label{tab:abl}
\begin{tabular}{llcccc}
\toprule
Model & Setting & test & train & change & Welch $p$ \\
\midrule
FNO & 4 modes & 0.0066 $\pm$ 0.0005 & 0.0048 $\pm$ 0.0001 & +36\% & 0.013 \\
FNO & 8 modes & 0.0048 $\pm$ 0.0001 & 0.0031 $\pm$ 0.0002 & $-$1\% & 0.897 \\
FNO & default (16 modes, $n$=400) & 0.0048 $\pm$ 0.0005 & 0.0032 $\pm$ 0.0001 & -- & -- \\
FNO & 32 modes & 0.0048 $\pm$ 0.0004 & 0.0029 $\pm$ 0.0002 & +0\% & 0.959 \\
FNO & no grid channel & 0.0043 $\pm$ 0.0005 & 0.0030 $\pm$ 0.0002 & $-$11\% & 0.279 \\
FNO & $n$=100 & 0.0159 $\pm$ 0.0016 & 0.0047 $\pm$ 0.0002 & +230\% & 0.004 \\
FNO & $n$=200 & 0.0087 $\pm$ 0.0004 & 0.0040 $\pm$ 0.0002 & +80\% & 0.001 \\
\midrule
DeepONet & $n$=100 & 0.3804 $\pm$ 0.0391 & 0.3581 $\pm$ 0.0489 & +75\% & 0.011 \\
DeepONet & $n$=200 & 0.2573 $\pm$ 0.0054 & 0.2363 $\pm$ 0.0104 & +18\% & 0.028 \\
DeepONet & default ($n$=400, 100 epochs) & 0.2173 $\pm$ 0.0144 & 0.1990 $\pm$ 0.0129 & -- & -- \\
DeepONet & 500 epochs & 0.1314 $\pm$ 0.0040 & 0.0833 $\pm$ 0.0030 & $-$40\% & 0.006 \\
DeepONet & 2000 epochs & 0.1276 $\pm$ 0.0081 & 0.0417 $\pm$ 0.0057 & $-$41\% & 0.002 \\
\midrule
MLP & default (100 epochs) & 0.0808 $\pm$ 0.0038 & 0.0232 $\pm$ 0.0006 & -- & -- \\
MLP & 500 epochs & 0.0766 $\pm$ 0.0050 & 0.0071 $\pm$ 0.0005 & $-$5\% & 0.314 \\
MLP & 2000 epochs & 0.0894 $\pm$ 0.0067 & 0.0307 $\pm$ 0.0155 & +11\% & 0.145 \\
\bottomrule
\end{tabular}
\end{table*}

\textbf{H3 (mode saturation): supported.} Going from 8 to 16 modes changes the mean error by about 1\%, below the 10\% criterion, whereas 4 modes is 36\% worse than 16 (Table~\ref{tab:abl}, Fig.~\ref{fig:sweeps}, left). The registered test compares means only and does not define ``clearly worse'' numerically; the Welch $p$ of the 4-mode row is given in Table~\ref{tab:abl} as a description, on three seeds. The 8-, 16- and 32-mode errors are indistinguishable at this sample size. The initial conditions have rapidly decaying spectrum amplitudes, so this saturation may be specific to the data distribution; we did not test another one.

\textbf{H4 (data efficiency): supported.} The FNO mean error is below the DeepONet mean error at 100, 200 and 400 pairs, by more than an order of magnitude at every size (Table~\ref{tab:abl}, Fig.~\ref{fig:sweeps}, right). Both models improve with more data; the FNO error falls about threefold from 100 to 400 pairs. Two caveats limit what this shows. First, the sweep fixes the number of epochs and the batch size, so a smaller training set also receives proportionally fewer gradient steps (500, 1000 and 2000 steps at 100, 200 and 400 pairs); data size and optimisation length are confounded and we ran no matched-step control. Second, DeepONet is under-fit at every size (training error close to test error), so its curve may reflect optimisation as much as data. H4 is a statement at a fixed 100-epoch budget, not a matched-convergence comparison.

\textbf{Training length for the cheap baselines.} Giving DeepONet 500 or 2000 epochs lowers its test error by about 40\%, yet it stays more than an order of magnitude above the FNO's (Table~\ref{tab:abl}). From 500 to 2000 epochs its test error barely moves while its training error halves, to about a third of the test error: with long training our DeepONet overfits the 400 pairs, so more epochs alone do not close the gap. The 100-epoch main comparison therefore understates DeepONet, but the ranking is unchanged. The MLP's test error at 500 and 2000 epochs is within noise of its 100-epoch value (three seeds, Welch $p$ of 0.314 and 0.145), and its training error is far below its test error at every length; we did not run a regularisation control. The CNN was not run for longer, because 500 epochs would exceed our per-run time limit (Table~\ref{tab:main}); whether it would improve is unknown.

\textbf{Grid channel.} Removing the coordinate channel does not hurt the FNO: the mean error is lower without it, but the difference is not distinguishable from seed noise on three seeds (Table~\ref{tab:abl}). On a periodic domain with a translation-equivariant operator the channel is expected to be unnecessary; this run is consistent with that expectation but does not test the mechanism.
